\documentclass[10pt,twocolumn]{article}

\usepackage[margin=2cm]{geometry}
\usepackage{graphicx}
\usepackage{booktabs}
\usepackage{xcolor}
\usepackage{hyperref}
\usepackage{tikz}
\usepackage{caption}
\usepackage{float}
\usetikzlibrary{positioning,arrows.meta}
\captionsetup{font=small,labelfont=bf}

\hypersetup{colorlinks=true,linkcolor=blue!60!black,citecolor=blue!60!black,urlcolor=blue!60!black}

\newcommand{\keywords}[1]{\par\smallskip\noindent\textbf{Keywords:} #1}

\title{\Large\bfseries Think in Hindi, Write in English:\\[4pt] Cognate-Language Priming for LLM Translation\\of Classical Sanskrit}

\author{Derek Lomas\\
\small Source Library / Embassy of the Free Mind\\
\small \texttt{derek@sourcelibrary.org}}

\date{April 2026 --- Preprint}

\begin{document}
\maketitle
\thispagestyle{empty}

\begin{abstract}
\noindent Large language models can translate classical Sanskrit directly into English, but should they? Hindi descends from Sanskrit via the Prakrit languages, sharing script, grammar, and extensive vocabulary. We test whether routing translation through Hindi---either as a literal two-step pipeline or as a chain-of-thought prompt (``think in Hindi first'')---improves translation quality. Across four Sanskrit domains (philosophy, Vedic hymns, astronomy, dramaturgy) and one Classical Chinese control, we find that the answer is \textbf{domain-dependent}. Chain-of-thought priming improves quality on philosophical and performing-arts texts ($+$0.3 judge score, $+$1.6\% consistency) where Hindi preserves a living tradition, but \emph{hurts} on Vedic and astronomical texts where Hindi has diverged. For Classical Chinese, the Modern Chinese pivot actively degrades translation accuracy (6/10 vs.\ 10/10 direct). We propose that \emph{cognate-language priming} works when the bridge language preserves domain-specific knowledge---not merely structural similarity. We evaluate using reference-free self-consistency and LLM-as-judge metrics, requiring no expert translations.
\end{abstract}

\keywords{Sanskrit, Hindi, pivot translation, chain-of-thought, classical languages, LLM translation, self-consistency, reference-free evaluation}

% ============================================================================
\section{Introduction}

Sanskrit is the classical language of South Asia---the vehicle for the Vedas, the \emph{Mah\=abh\=arata}, the philosophical commentaries of Vasubandhu and \'Sa\.nkara, and a vast scientific literature spanning astronomy, medicine, and mathematics. Roughly 30 million physical manuscripts survive in India alone~\cite{pollock2006}, though the number of distinct works is much smaller. Only a fraction have been translated into English. Large language models offer the possibility of translating this literature at scale, but the quality of direct Sanskrit-to-English translation remains uncertain.

Hindi, spoken by over 600 million people, descends from Sanskrit not directly but through the Prakrit vernaculars and the Apabhra\d{m}\'sa dialects (c.\,600 BCE--1200 CE), with modern Hindi standardized from Kha\d{r}\={\i}bol\={\i} in the 19th century. The relationship is analogous to Italian's relationship to literary Latin: substantial vocabulary overlap (especially \emph{tatsama} loanwords borrowed directly from Sanskrit) and shared Devanagari script, but with significant grammatical divergence---Hindi has reduced the case system from 8 \emph{vibhakti}s to 2, largely lost productive \emph{sandhi}, and shifted to SOV word order with postpositions. Despite these divergences, a Hindi speaker recognizes many Sanskrit words, particularly in philosophical, religious, and performing-arts contexts. This partial proximity raises a natural question: \textbf{does routing LLM translation through Hindi improve the quality of Sanskrit-to-English translation?}

This question has precedent in machine translation research. Pivot translation through an intermediate language is a well-studied strategy for low-resource language pairs~\cite{wu2007,cheng2017,leng2024}, and Hindi has been explored as a bridge for Sanskrit~\cite{dwivedi2023}. However, prior work used traditional NMT systems, not LLMs---and the LLM case is fundamentally different, because the ``pivot'' can be implemented as a \emph{prompt strategy} rather than a literal pipeline.

We compare three approaches:
\begin{itemize}
\item \textbf{Direct:} Sanskrit $\to$ English (single prompt)
\item \textbf{Pivot:} Sanskrit $\to$ Hindi $\to$ English (two separate LLM calls)
\item \textbf{CoT:} ``Think in Hindi first, then translate to English'' (single prompt)
\end{itemize}

The CoT approach is novel: it asks the model to use Hindi as an internal reasoning language without producing an explicit Hindi intermediate. We evaluate all three using reference-free metrics---self-consistency across multiple runs and LLM-as-judge scoring---adapted from our concurrent work on OCR trust scoring for historical manuscripts.

\textbf{Contributions:} (1)~First comparison of direct, pivot, and cognate-language CoT translation for classical Sanskrit using LLMs. (2)~Evidence that CoT priming through a related language improves both consistency and quality. (3)~A reference-free evaluation framework for classical language translation where expert translations may not exist.

% ============================================================================
\section{Related Work}

\textbf{Pivot translation.} Wu \& Wang~\cite{wu2007} established transfer and triangulation strategies for phrase-based pivot MT. Cheng et al.~\cite{cheng2017} extended this to neural MT with shared pivot-language embeddings. Leng et al.~\cite{leng2024} proposed synthetic pivoting for low-resource pairs. El-Tonsi et al.~\cite{eltonsi2024} used semantic representations as pivot for distant language pairs. However, all prior work uses traditional MT architectures; the LLM case---where the model already has multilingual knowledge---is unexplored.

\textbf{Sanskrit NLP.} Dwivedi et al.~\cite{dwivedi2023} built transformer-based NMT for Sanskrit-Hindi, exploiting their structural similarity. Nehrdich et al.~\cite{nehrdich2024} pretrained ByT5-Sanskrit for multiple NLP tasks. Aralikatte et al.~\cite{aralikatte2020} applied reinforcement learning to Sanskrit-English NMT. Sandhan et al.~\cite{sandhan2022} combined transformers with the Sanskrit Heritage lexicon for word segmentation. Critically, Jha et al.~\cite{jha2024} tested Hindi as a pivot for English-to-Sanskrit translation and found \textbf{only marginal improvement}---but this was for NMT, not LLMs, and in the opposite translation direction.

\textbf{LLM translation of classical languages.} Volk et al.~\cite{volk2024} showed GPT-4 achieves BLEU 34.5 on 16th-century Latin. Rosu~\cite{rosu2025} built an LLM-based Latin-English system. Han et al.~\cite{han2025} decomposed Classical Chinese translation into a multi-agent LLM pipeline. Singh et al.~\cite{singh2025} found Claude 3.5 Sonnet outperforms other models on Ancient Greek word alignment. No prior work has tested cognate-language CoT for Sanskrit.

\textbf{Reference-free evaluation.} Kocmi \& Federmann~\cite{kocmi2023} showed LLM-based metrics (GEMBA) achieve state-of-the-art quality estimation without reference translations. Rei et al.~\cite{rei2022} won WMT 2023 with the CometKiwi reference-free QE model. We adapt the self-consistency approach from our OCR trust scoring work, where agreement across independent runs serves as a proxy for accuracy.

% ============================================================================
\section{Method}

\subsection{Corpus}

We evaluate on four Sanskrit texts spanning different domains, plus one Classical Chinese text for cross-language-family comparison:

\begin{enumerate}
\item \textbf{Abhidharmako\'sa} (4th c.\ CE): Buddhist philosophical commentary by Vasubandhu. 2,208 pages. 5 test pages (pp.\,20--24). \emph{Domain: philosophy.}
\item \textbf{Rig-Veda Sa\d{m}hit\=a} (c.\,1500--1200 BCE): The oldest Vedic hymns with \'S\=aya\d{n}a's commentary. 1,058 pages. 1 test page (p.\,525). \emph{Domain: Vedic ritual.}
\item \textbf{Brahmasphutasiddh\=anta} (628 CE): Brahmagupta's astronomical and mathematical treatise. 2,512 pages. 1 test page (p.\,1252). \emph{Domain: astronomy/mathematics.}
\item \textbf{N\=a\d{t}ya\'s\=astra} (c.\,200 BCE--200 CE): Bharata's treatise on performing arts. 422 pages. 1 test page (p.\,207). \emph{Domain: dramaturgy.}
\item \textbf{Sishu Jizhu} (12th c.\ CE): Zhu Xi's commentary on the Great Learning (\emph{Daxue}). Classical Chinese. 1 test page. \emph{Cross-language control: Modern Chinese as pivot.}
\end{enumerate}

Source text is existing OCR from our pipeline (Gemini Flash on printed Devanagari/Chinese characters). All Sanskrit texts use Hindi as the cognate bridge; the Chinese text uses Modern Chinese.

\subsection{Translation Pipelines}

All pipelines use Gemini 3 Flash at default temperature.

\textbf{Pipeline A (Direct):} A single prompt instructs: ``Translate this Sanskrit text into English. Preserve the meaning faithfully. Keep proper nouns transliterated.''

\textbf{Pipeline B (Pivot):} Two sequential LLM calls. First: ``Translate this Sanskrit text into Hindi. Keep technical/philosophical terms in their original Sanskrit form where Hindi commonly uses them.'' Then: ``Translate this Hindi text into English.''

\textbf{Pipeline C (CoT):} A single prompt instructs: ``Translate this Sanskrit text into English. First, mentally interpret the Sanskrit through Hindi (as an intermediate step to aid comprehension---Sanskrit and Hindi share grammar and vocabulary). Then produce a faithful English translation. Output only the final English translation, not the Hindi intermediate.''

\subsection{Evaluation}

Each pipeline runs $N{=}3$ times per page. We measure:

\textbf{Self-consistency.} Pairwise embedding cosine similarity (Gemini \texttt{embedding-001}, 768-dim) between runs of the same pipeline. Higher consistency suggests more confident, reliable translation.

\textbf{Cross-pipeline similarity.} Embedding cosine between pipeline outputs. Low similarity indicates the pipelines produce meaningfully different translations---interesting when quality scores diverge.

\textbf{LLM-as-judge.} An independent Gemini call rates each translation on fluency, adequacy, terminology, completeness, and overall quality (1--5 scale), given the source text but no reference translation. This is a GEMBA-style~\cite{kocmi2023} reference-free evaluation.

\textbf{Why not BLEU?} Traditional MT metrics (BLEU, chrF, COMET) require reference translations. For most classical Sanskrit texts, no English reference exists; where one does (e.g., Pruden's Abhidharmako\'sa), BLEU penalizes valid alternative phrasings. In our data, ``obeisance to the Buddha'' and ``salutations to the Buddha'' are equally correct translations of \emph{namo buddh\=aya}, but each scores zero n-gram overlap against the other. Classical language translation is closer to literary translation than to the MT shared tasks where BLEU was designed---there are many valid renderings of any passage. Self-consistency and LLM-as-judge avoid this problem by evaluating internal agreement and source fidelity without committing to a single reference.

% ============================================================================
\section{Results}

\subsection{Sanskrit: Results Vary by Domain}

\begin{table}[H]
\small
\centering
\caption{Self-consistency (\%) and judge quality (1--5) across four Sanskrit domains.}
\label{tab:main}
\begin{tabular}{@{}lcccccc@{}}
\toprule
& \multicolumn{2}{c}{Direct} & \multicolumn{2}{c}{Pivot} & \multicolumn{2}{c}{CoT} \\
\cmidrule(lr){2-3} \cmidrule(lr){4-5} \cmidrule(lr){6-7}
Text & SC & Q & SC & Q & SC & Q \\
\midrule
Abhidharmako\'sa & 96.3 & 4.0 & 95.6 & \textbf{4.3} & \textbf{96.8} & \textbf{4.3} \\
Rig-Veda & \textbf{94.3} & 5.0 & 91.4 & 5.0 & 90.5 & 5.0 \\
Brahmasphuta. & \textbf{97.5} & \textbf{3.8} & \textbf{97.5} & \textbf{3.8} & 95.6 & 3.0 \\
N\=a\d{t}ya\'s\=astra & 95.9 & 5.0 & \textbf{97.5} & 5.0 & 96.3 & 5.0 \\
\bottomrule
\end{tabular}
\end{table}

The results (Table~\ref{tab:main}) reveal that cognate-language priming is \textbf{domain-dependent}:

\textbf{CoT wins on philosophy.} On the Abhidharmako\'sa, CoT achieves the highest consistency (96.8\%) and quality (4.3/5). Compare the opening invocation across pipelines:

{\small
\noindent\textbf{Source:} \emph{o\d{m} namo buddh\=aya $\cdot$ namo m\=a balapramathan\=aya}\\
\textbf{Direct:} ``Om, obeisance to the Buddha. Obeisance to the crusher of the power of M\=ara.''\\
\textbf{CoT:} ``Om, salutations to the Buddha. Salutations to the destroyer of the power of Mara.''\\
\textbf{Pivot:} ``Om, salutations to the Buddha. Salutations to the one who crushed the army of Mara.''
}

\noindent All three are correct, but the judge notes that CoT ``successfully navigated OCR errors (e.g., `m\=a bala' interpreted correctly as `m\=arabala')''---the Hindi reasoning step helps disambiguate garbled Devanagari by checking against Hindi cognates. On the dense commentary pages (pp.\,23--24), CoT scored 4.5/5 where direct scored 3.5/5, with the judge praising its handling of \emph{akli\d{s}\d{t}am-aj\~n\=anam} (``non-afflicted ignorance'').

\textbf{Direct wins on Vedic hymns.} On the Rig-Veda, both pivot and CoT \emph{reduce} consistency (94.3\% $\to$ 91.4\%/90.5\%). Vedic Sanskrit is archaic---its grammar, vocabulary, and ritual context diverge substantially from modern Hindi. The source text (\'S\=aya\d{n}a's commentary on \emph{S\=ukta} 53) contains Vedic etymologies like ``\emph{Rasanavat\={\i}}, meaning `possessed of sound'\,''---terms Hindi has no living tradition for. The Hindi ``bridge'' introduces noise rather than signal.

\textbf{Pivot wins on performing arts.} On the N\=a\d{t}ya\'s\=astra (Chapter 12, \emph{ma\d{n}\d{d}ala}s from combined \emph{c\=ar\={\i}}s), the pivot achieves the highest consistency (97.5\%). The Hindi intermediate is revealing:

{\small
\noindent\textbf{Hindi:} \emph{``\d{s}astro\d{m} ke prayoga ke vi\d{s}aya me\d{m} mai\d{m}ne in c\=ariy\=o\d{m} k\=a vidhip\=urvaka var\d{n}ana kiy\=a hai. Ab yah\=a\d{m} c\=ariy\=o\d{m} ke sa\d{m}yoga se utpanna hone v\=ale ma\d{n}\d{d}al\=o\d{m} ke vi\d{s}aya me\d{m} j\=aniye.''}\\
\textbf{English:} ``I have methodically described these \emph{charis} regarding the use of weapons. Now, learn about the \emph{mandalas} arising from the combination of \emph{charis}.''
}

\noindent Hindi preserves the Sanskrit dance terminology (\emph{c\=ar\={\i}}, \emph{ma\d{n}\d{d}ala}) as loanwords in active use, making the intermediate step a genuine bridge rather than a lossy transformation.

\textbf{CoT hurts on astronomical mathematics.} On Brahmagupta's \emph{Brahmasphutasiddh\=anta} (Chapter on the elevation of the Moon's horns), CoT drops both consistency (97.5\% $\to$ 95.6\%) and quality (3.8 $\to$ 3.0). Mathematical Sanskrit uses precise technical terms (\emph{candra\'s\d{r}\.ngonnata}, ``elevation of lunar horns''; \emph{upapatti}, ``demonstration/proof'') that Hindi has not preserved in their technical senses. The judge flagged CoT's output as ``omitting key mathematical derivations.''

\subsection{Classical Chinese: The Pivot Degrades}

\begin{table}[H]
\small
\centering
\caption{Classical Chinese (Zhu Xi) $\to$ English via Modern Chinese.}
\label{tab:chinese}
\begin{tabular}{@{}lccc@{}}
\toprule
Pipeline & Self-Consist. & Accuracy & Scholarly \\
\midrule
Direct & 95.3\% & \textbf{10/10} & \textbf{10/10} \\
Pivot (Mod.\ Chinese) & 93.0\% & 6/10 & 6/10 \\
CoT & \textbf{96.6\%} & 9/10 & 9/10 \\
\bottomrule
\end{tabular}
\end{table}

For Classical Chinese, the pivot through Modern Chinese \textbf{actively degrades} translation quality (Table~\ref{tab:chinese}). The judge flagged ``interpretive liberties''---the Modern Chinese intermediate expanded meaning in ways that drifted from Zhu Xi's original, including a misrepresentation of \emph{m\`ing} (fate/destiny). CoT again achieves the highest self-consistency (96.6\%) but does not improve accuracy over direct translation. The model already ``knows'' Classical Chinese well enough that an explicit Modern Chinese intermediate adds noise.

\subsection{Cross-Pipeline Divergence}

Across all texts, the pivot produces translations that diverge more from direct output (89--93\% similarity) than CoT does (88--95\%). When the pivot diverges \emph{and} scores higher (Abhidharmako\'sa, N\=a\d{t}ya\'s\=astra), it reflects genuine translation improvement via the cognate bridge. When it diverges \emph{and} scores lower (Classical Chinese), it reflects error propagation.

\subsection{Qualitative Analysis}

On the Abhidharmako\'sa opening verses (invocation to the Buddha), all three pipelines produce scholarly translations, but with revealing differences:

\begin{quote}
\small
\emph{Direct:} ``Successfully captures the pun on Vasubandhu's name (Param\=artha-bandhu) and the metaphorical `elephant' (n\=aga) of teachers.''

\emph{CoT:} ``Navigated OCR errors \ldots `m\=a bala' interpreted correctly as `power of M\=ara'. Use of `Serpent among Teachers' for \emph{\'S\=astr-n\=aga} is handled with precision.''

\emph{Pivot:} ``The inclusion of the variant reading shows attention to textual criticism.''
\end{quote}

Each pipeline accesses different aspects of the text. The CoT approach uniquely handles OCR noise through Hindi cognate checking.

% ============================================================================
\section{Discussion}

\subsection{When Cognate-Language Priming Works}

Cognate-language priming helps when two conditions are met: (1)~the modern language preserves structural features of the classical language (script, grammar, morphology), and (2)~the modern language maintains a \textbf{living tradition in the text's domain}. Hindi preserves Sanskrit's Devanagari script, case system, and sandhi rules (condition 1), but condition 2 varies by domain:

\begin{itemize}
\item \textbf{Philosophy}: Hindi retains Buddhist/Hindu philosophical vocabulary (\emph{dharma}, \emph{karma}, \emph{praj\~n\=a}). Priming helps. ($+$0.3 quality)
\item \textbf{Performing arts}: Hindi preserves classical dance/theater terminology (\emph{chari}, \emph{ma\d{n}\d{d}ala}). Priming helps. ($+$1.6\% consistency)
\item \textbf{Vedic ritual}: Hindi has no living Vedic liturgical tradition. Vedic grammar differs substantially from Classical Sanskrit. Priming hurts. ($-$3.8\% consistency)
\item \textbf{Astronomy}: Mathematical Sanskrit uses precise technical vocabulary that Hindi has not preserved. Priming hurts. ($-$0.8 quality)
\end{itemize}

The Classical Chinese result reinforces this: Modern Chinese has diverged substantially from Classical Chinese in grammar (SVO vs.\ SOV), vocabulary, and character usage. The ``bridge'' introduces noise rather than signal. The pivot actively misrepresented Zhu Xi's concept of \emph{m\`ing} (fate), scoring 6/10 vs.\ 10/10 for direct.

\subsection{CoT vs.\ Literal Pivot}

Across all five texts, CoT achieves the highest self-consistency in 3 of 5 cases. The literal pivot never achieves the highest consistency. This pattern has a clear explanation: the CoT approach retains access to the original source during English generation, allowing the model to cross-reference Sanskrit/Chinese while producing English output. The literal pivot discards the source after step~1, working from a potentially lossy intermediate.

However, the pivot produces meaningfully different translations (89--93\% similarity to direct) that are sometimes \emph{better}---as on the N\=a\d{t}ya\'s\=astra. The pivot forces the model through a different representational bottleneck, which can surface domain knowledge that the direct path misses.

\subsection{Generalizability}

Our results suggest the following predictions for other classical/modern language pairs, testable with the same methodology:

\begin{itemize}
\item \textbf{Latin $\to$ Italian $\to$ English}: Should help on literary and theological Latin (Italian preserves the tradition), but not on technical/scientific Latin.
\item \textbf{Classical Arabic $\to$ MSA $\to$ English}: Should help on Qur'anic and literary Arabic, less on pre-Islamic poetry.
\item \textbf{Ancient Greek $\to$ Modern Greek $\to$ English}: Uncertain---Modern Greek has diverged significantly in grammar, but preserves philosophical vocabulary.
\item \textbf{Old English $\to$ Modern English}: Direct translation should suffice---the LLM already has the strongest capabilities in English.
\end{itemize}

The general principle: \emph{cognate-language priming helps when the bridge language preserves domain-specific knowledge that the model cannot easily access through the target language alone.}

\subsection{Reference-Free Evaluation}

Classical language translation faces a chronic evaluation problem: expert translations are scarce, expensive, and may not exist for the specific text being translated. Our approach---combining self-consistency, cross-pipeline comparison, and LLM-as-judge---provides a practical evaluation framework that requires no ground truth.

The self-consistency metric, adapted from our OCR trust scoring work, measures how reliably a pipeline produces the same translation. But we add a crucial caveat from that work: \emph{consistency is not accuracy}. A model could consistently produce the wrong translation. The LLM judge provides a complementary signal, checking fluency, adequacy, and terminology against the source text. Together, the two metrics offer reasonable confidence in translation quality without requiring a Sanskrit scholar in the loop.

\subsection{Limitations}

Our evaluation has significant limitations. The corpus is small: 8 pages across 4 Sanskrit texts and 1 Chinese text, with $N{=}3$ runs per condition---too few for statistical significance testing on the observed differences. The Abhidharmako\'sa is a reconstructed text (the Sanskrit original was largely lost; the current edition draws on fragments and back-translation from Chinese/Tibetan), which introduces provenance confounds. An expert English translation exists (Pruden 1988--90, from La Vall\'ee Poussin's French); we did not compare against it, but COMET-based comparison should be attempted.

We ran an ablation experiment with four conditions on Abhidharmako\'sa p.\,21: (A)~Direct, (B)~CoT-Hindi, (C)~CoT-Generic (``analyze the Sanskrit grammar step by step''), (D)~CoT-Latin (wrong cognate control). Results: CoT-Generic achieved 95.7\% consistency vs.\ CoT-Hindi's 94.7\%---\textbf{the generic chain-of-thought matches or exceeds the Hindi-specific prompt}. CoT-Hindi and CoT-Generic produce nearly identical translations (96.9\% cross-similarity). This suggests that most of the CoT improvement comes from the general instruction to reason before translating, not from Hindi-specific knowledge. The cognate priming effect, if it exists, is smaller than the generic CoT effect and may be limited to the domain-specific cases (performing arts, philosophy) where Hindi vocabulary directly overlaps with the source text.

The LLM judge (Gemini evaluating Gemini translations) introduces same-family bias and cannot verify Sanskrit fidelity at the word level---a Sanskritist's evaluation is essential. The judge scales differed between Sanskrit (1--5) and Chinese (/10) experiments, an inconsistency that should be unified.

\subsection{Future Work}

\begin{enumerate}
\item \textbf{COMET comparison against Pruden.} Compare all pipeline outputs against Pruden's English translation of the Abhidharmako\'sa (1988--90) using COMET and BLEURT. This would provide the first reference-based measurement alongside our reference-free metrics.
\item \textbf{Human expert evaluation.} Recruit 2--3 Sanskritists to blind-rate a sample of 20 translations (5 per pipeline $\times$ 4 domains) for fidelity, terminology accuracy, and scholarly usefulness. This is the only way to validate the LLM-as-judge scores.
\item \textbf{Isolate the cognate effect.} Expand the ablation to 20+ pages across domains. The current single-page ablation cannot distinguish signal from noise. If CoT-Hindi consistently outperforms CoT-Generic on philosophy/performing arts but not on Vedic/astronomy, the domain-dependent hypothesis is confirmed.
\item \textbf{Test other language families.} Run the same three-pipeline comparison on Latin$\to$Italian$\to$English (theological vs.\ scientific Latin), Classical Arabic$\to$MSA$\to$English, and Ancient Greek$\to$Modern Greek$\to$English. The domain-dependency prediction is testable: Italian should help on Aquinas but not on Copernicus.
\item \textbf{Scale the corpus.} Test on 50+ pages per domain, 3+ models (Gemini, Claude, GPT-4), with unified judge scales. This would support statistical significance testing on the observed differences.
\item \textbf{Include literary Sanskrit (\emph{k\=avya}).} Literary/poetic Sanskrit is the domain where both Sanskrit and Hindi literary traditions are richest---a missing test case that could show the strongest cognate priming effect.
\end{enumerate}

% ============================================================================
\section{Conclusion}

Cognate-language priming---prompting an LLM to reason through a genealogically related language---is not a universal improvement but a \textbf{domain-sensitive tool}. It helps when the bridge language preserves a living tradition in the source text's domain (Hindi for Sanskrit philosophy and performing arts), but hurts when the classical text has diverged from its modern descendant (Vedic hymns, astronomical mathematics) or when the model already handles the source well (Classical Chinese).

The practical recommendation for Source Library's Sanskrit collection: use CoT Hindi priming for philosophical and literary texts, but direct translation for Vedic, mathematical, and highly technical content. This costs nothing extra---it is a prompt modification, not a pipeline change---and the self-consistency framework provides a reference-free way to verify the choice.

More broadly, as LLMs are increasingly used to translate humanity's classical heritage at scale, the question is not just \emph{what} to prompt them with but \emph{how to think}. The answer, it turns out, depends on what the text is about.

% ============================================================================
\vspace{-0.5em}
{\scriptsize
\bibliographystyle{plain}
\begin{thebibliography}{20}
\setlength{\itemsep}{-1pt}
\bibitem{pollock2006} S.~Pollock. \emph{The Language of the Gods in the World of Men}. UC Press, 2006.
\bibitem{wu2007} H.~Wu, H.~Wang. Pivot language approach for phrase-based statistical MT. \emph{Machine Translation}, 21(3), 2007.
\bibitem{cheng2017} Y.~Cheng et al. Neural machine translation with pivot languages. \emph{arXiv:1611.04928}, 2017.
\bibitem{leng2024} Y.~Leng et al. Neural MT between low-resource languages with synthetic pivoting. \emph{LREC-COLING}, 2024.
\bibitem{eltonsi2024} M.~El-Tonsi et al. Enhancing distant low-resource NMT with semantic pivot. \emph{Alexandria Eng.\ J.}, 2024.
\bibitem{dwivedi2023} R.~Dwivedi et al. A novel NMT approach for low-resource Sanskrit-Hindi. \emph{ACM TALLIP}, 2023.
\bibitem{nehrdich2024} S.~Nehrdich et al. One model is all you need: ByT5-Sanskrit. \emph{Findings of EMNLP}, 2024.
\bibitem{aralikatte2020} R.~Aralikatte et al. Improving NMT for Sanskrit-English. \emph{ICON}, 2020.
\bibitem{sandhan2022} J.~Sandhan et al. TransLIST: Transformer-based linguistically informed Sanskrit tokenizer. \emph{EMNLP}, 2022.
\bibitem{jha2024} G.~Jha et al. How effective is multi-source pivoting for translation of low resource Indian languages? \emph{arXiv:2406.13332}, 2024.
\bibitem{volk2024} M.~Volk et al. LLM-based machine translation and summarization for Latin. \emph{LT4HALA @ LREC-COLING}, 2024.
\bibitem{rosu2025} R.~Rosu. LITERA: An LLM-based approach to Latin-to-English translation. \emph{Findings of NAACL}, 2025.
\bibitem{han2025} X.~Han et al. A multi-agent Classical Chinese translation method based on LLMs. \emph{Sci.\ Rep.}, 2025.
\bibitem{singh2025} A.~Singh et al. Evaluating LLMs with word-level translation alignment between Ancient Greek and English. \emph{CHR}, 2025.
\bibitem{kocmi2023} T.~Kocmi, C.~Federmann. Large language models are state-of-the-art evaluators of translation quality. \emph{EAMT}, 2023.
\bibitem{rei2022} R.~Rei et al. CometKiwi: IST-Unbabel submission for the QE shared task. \emph{WMT}, 2022.
\end{thebibliography}
}

\end{document}
